% AUTHOR: Amlal El Mahrouss
% PURPOSE: M3.

\documentclass[10pt]{article}

%% ---- Core Packages ----------------------------------------
\usepackage[T1]{fontenc}
\usepackage[utf8]{inputenc}
\usepackage{lmodern}
\usepackage{microtype}

%% ---- Page Geometry ----------------------------------------
\usepackage[
top=1in, bottom=1in,
left=0.75in, right=0.75in,
columnsep=0.25in
]{geometry}

%% ---- Mathematics ------------------------------------------
\usepackage{amsmath, amssymb, amsthm, amsfonts}
\usepackage{mathtools}
\usepackage{bm}
\usepackage{textcomp}
\usepackage{esint}

%% ---- Figures & Tables -------------------------------------
\usepackage{graphicx}
\usepackage{booktabs}
\usepackage{multirow}
\usepackage{array}
\usepackage{caption}
\usepackage{subcaption}
\usepackage{float}

%% ---- Code Listings ----------------------------------------
\usepackage{listings}
\usepackage{xcolor}
\usepackage{algorithm}
\usepackage{algpseudocode}

\lstdefinestyle{codestyle}{
	backgroundcolor=\color{gray!8},
	basicstyle=\ttfamily\footnotesize,
	breakatwhitespace=false,
	breaklines=true,
	captionpos=b,
	commentstyle=\color{green!50!black},
	keywordstyle=\color{blue}\bfseries,
	stringstyle=\color{orange!70!black},
	numberstyle=\tiny\color{gray},
	numbers=left,
	numbersep=5pt,
	showstringspaces=false,
	frame=single,
	rulecolor=\color{gray!40},
	tabsize=2
}
\lstset{style=codestyle}

%% ---- Hyperlinks -------------------------------------------
\usepackage[
colorlinks=true,
linkcolor=blue!70!black,
citecolor=green!50!black,
urlcolor=blue!60!black
]{hyperref}
\usepackage{url}

%% ---- Bibliography -----------------------------------------
\usepackage[numbers, sort&compress]{natbib}

%% ---- Miscellaneous ----------------------------------------
\usepackage{enumitem}
\usepackage{siunitx}
\usepackage{cleveref}

%% ---- Theorem Environments ---------------------------------
\theoremstyle{plain}
\newtheorem{theorem}{Theorem}[section]
\newtheorem{lemma}[theorem]{Lemma}
\newtheorem{corollary}[theorem]{Corollary}
\newtheorem{proposition}[theorem]{Proposition}

\theoremstyle{definition}
\newtheorem{definition}[theorem]{Definition}
\newtheorem{example}[theorem]{Example}

\theoremstyle{remark}
\newtheorem{remark}[theorem]{Remark}

%% ---- Custom Commands --------------------------------------
\newcommand{\R}{\mathbb{R}}
\newcommand{\N}{\mathbb{N}}
\newcommand{\Z}{\mathbb{Z}}
\newcommand{\E}{\mathbb{E}}
\newcommand{\Prob}{\mathbb{P}}
\newcommand{\norm}[1]{\left\lVert #1 \right\rVert}
\newcommand{\abs}[1]{\left\lvert #1 \right\rvert}
\newcommand{\inner}[2]{\left\langle #1,\, #2 \right\rangle}
\newcommand{\ie}{\textit{i.e.}\xspace}
\newcommand{\eg}{\textit{e.g.}\xspace}
\newcommand{\etal}{\textit{et al.}\xspace}

%% ---- Caption Formatting -----------------------------------
\captionsetup{
	font=small,
	labelfont=bf,
	format=hang,
	justification=justified
}

\renewcommand{\qedsymbol}{\ensuremath{\blacksquare}}

\usepackage{lettrine}

%% ---- Header / Footer --------------------------------------
\usepackage{fancyhdr}
\pagestyle{fancy}
\fancyhf{}
\fancyhead[R]{\small\thepage}
\renewcommand{\headrulewidth}{0.4pt}

%% ============================================================
%%  FRONT MATTER
%% ============================================================

\title{%
	\vspace{-1.5em}%
	\large\textbf{Second Additional Write-up for Future Problems.}%
}
\author{By Amlal El Mahrouss}
\date{\today}

%% ============================================================
\begin{document}
%% ============================================================

\maketitle
\tableofcontents
\newpage

%% ============================================================
\section{Abstract:}
%% ============================================================

%% ============================================================
\section{Lemmas and Proofs:}
%% ============================================================

\lettrine[lines=1, lhang=0.1, loversize=0.2]{T}HE following notes are stated in a reusable manner for future work in Abstract Algebra and Differential Geometry.

%% ============================================================
\subsection{Lemmas and Theorems for Tensors:}
%% ============================================================

\begin{lemma}
\label{lem:1}
There exist a following identity of factor $p$ of an index of a tensor $\mathcal{I}_{kij}$ at $i, j, k$ in $\mathbb{N}^{+}$.
\begin{equation}
	||\mathcal{E}_{ijk}||^{p} = (p \cdot \mathcal{I}_{kij}).
\end{equation}
The following equation holds for $\forall \mathcal{E} \neq 0 \land \forall p \neq 0.$\\
For $i, j, k \subset \mathbb{N}^{+}$.
\end{lemma}
\begin{remark}
Per Lemma~\ref{lem:1} There exist a specific nonlinear graph of $\mathcal{E}$, if and only if $p < 0$.
\end{remark}

\begin{definition}
\label{def:1}
Per \ref{lem:1}. Let the equation of $p, n \in \mathbb{N}^{+}$ be defined for two tensors $\mathcal{E}_{ijk} \land \mathcal{I}_{kij}$.
\begin{equation}
	||\mathcal{E}_{ijk}||^{p+n} = (p \cdot \mathcal{I}_{kij})^{n}.
\end{equation}
\begin{remark}
\label{rem:1}
We could divide to $p$ on both sides of the equation per \ref{def:1}:
\begin{align*}
	||\mathcal{E}_{ijk}||^{p+n} = (p \cdot \mathcal{I}_{kij})^{n},\\
	\frac{||\mathcal{E}_{ijk}||^{p+n}}{p} = \mathcal{I}_{kij}^{n}, \, \forall \mathcal{E}_{ijk} \neq 0, \, \forall p \neq 0, \, \forall \mathcal{I}_{kij}^n \neq 0, \, n \geq 1.
\end{align*}
Such properties hold for $i, j, k \subset \mathbb{N}^{+}$.
\end{remark}
\end{definition}
\begin{remark}
	Per \ref{rem:1} $i \land j \land k$ shall be non-zero values in $\mathbb{N}$. Moreover the identity shall hold for: $\mathcal{I}_{kij}^{n} \neq \mathcal{E}_{ijk}^{n}$.
\end{remark}

\begin{theorem}
There exist a sequence of products of manifold tensors $\mathcal{T}$ expressed as $\mathcal{B}$ indexed $i \in \mathbb{N}^{+}$:
\begin{equation}
	\mathcal{B}_{i}(n) = \prod_{k=1}^{n} \mathcal{T}_{k}^{i} + C_{i}, \, n \in \mathbb{N}^{+}.
\end{equation}
Where $\mathcal{T}$ of index $k$ is a tensor representing a manifold of four diemensional properties. And $C$ is a residue constant in a arbitary n-plane. Moreover the inverse function of $\mathcal{B}_{i}(n)$ is:
\begin{equation}
	\mathcal{B}_{i}(n)^{-1} = \pm \frac{\mathcal{I}^{-1}}{\prod_{k=1}^{n} \mathcal{T}_{k}^{i} + C_{i}}, \, n, i \in \mathbb{N}^{+}.
\end{equation}
Where $\mathcal{I}^{-1}$ is an non-zero inverse tensor.
\end{theorem}

\section{Theorems, Definitions, and Several PDEs:}
The following definition applies to Machine Learning for specific training techniques:

\begin{definition}
\label{def:eq-tensor}
Consider the following function $\mathcal{G}$ being equal to:
\begin{equation}
	\mathcal{G}(x) = {\mathcal{T}}\Bigl(\Bigl( \int_{n}^{p} Q(x) \, dx \Bigr) + \mathcal{D}_t\Bigr).
\end{equation}
And the tensor $\mathcal{Z}$ being equal to the following Ricci tensor at $t \in \mathbb{N}^{+}$:
\begin{equation}
	\mathcal{Z} = \frac{1}{4} \nabla \mathcal{R}_{t}.
\end{equation}
Let following equation be in a linear 3-manifold space in $\mathbb{R}$:
\begin{equation}
	\Delta \mathcal{X} = \sum_{n=1}^{\infty} \sum_{t=1}^{\infty} {\mathcal{T}}\Bigl(\Bigl( \int_{n}^{p} Q(x) \, dx \Bigr) + \mathcal{D}_t\Bigr) \pm \frac{1}{4} \nabla \mathcal{R}_{t}.
\end{equation}
\end{definition}

\subsection{Severel theorems per \ref{def:eq-tensor}:}

%% ============================================================
\section{Conclusion:}
%% ============================================================
Finishing of per previous definitions (\ref{def:1}, \ref{lem:1}), this note explores several identities and proofs regarding tensors.

\nocite{*}
\bibliographystyle{acm}
\bibliography{pinter2010book, weil1967basic}

%% ============================================================
\end{document}
%% ============================================================
